% Supplementary material for the Journal of Industrial Ecology cut
% Extended theorem inventory, delay-dynamics interface, full application tables,
% and an operational weak/strong-sustainability formalization.
\documentclass[11pt]{article}
\usepackage[a4paper,margin=1in]{geometry}
\usepackage[T1]{fontenc}
\usepackage{lmodern}
\usepackage{amsmath,amssymb,mathtools}
\usepackage{booktabs,longtable,tabularx,array}
\usepackage{enumitem}
\usepackage{xcolor}
\usepackage{xurl}
\usepackage[colorlinks=true,allcolors=blue!45!black]{hyperref}
\usepackage{microtype}
\setlength{\parindent}{0pt}
\setlength{\parskip}{5pt}
\setlength{\emergencystretch}{3em}
\sloppy
\setlist[itemize]{leftmargin=1.5em,itemsep=2pt,topsep=3pt}
\setlist[enumerate]{leftmargin=1.7em,itemsep=3pt,topsep=3pt}
\setlength{\tabcolsep}{4pt}
\renewcommand{\arraystretch}{1.18}
\newcolumntype{Y}{>{\raggedright\arraybackslash}X}
\newcolumntype{P}[1]{>{\raggedright\arraybackslash}p{#1}}
\newcommand{\est}{\texttt{established}}
\newcommand{\nesta}{\texttt{not established}}
\newcommand{\na}{\texttt{not applicable}}
\newcommand{\R}{\mathbb{R}}
\newcommand{\Rpos}{\mathbb{R}_{+}}
\newcommand{\dd}{\,\mathrm{d}}

\begin{document}

\title{Supplementary material\\
\large Typed material-flow ledgers for componentwise sustainability diagnostics}
\author{Amin Abaee\\[0.25em]
{\small Independent Researcher, Tehran, Iran}\\[0.25em]
{\small\href{https://orcid.org/0000-0002-0019-1842}{\nolinkurl{ORCID: 0000-0002-0019-1842}}}\\[0.25em]
{\small\href{mailto:amin_abaee@ut.ac.ir}{\nolinkurl{amin_abaee@ut.ac.ir}}}}
\date{19 September 2026}
\maketitle

\begin{abstract}
The theorem inventory, delay-dynamics interface, application tables and operational weak- and strong-sustainability regimes specify the typed material-flow ledger. The formalization is conditional on declared types, routes, barriers, demand and horizon; it does not constitute a universal theorem about sustainability ethics or an empirical calibration of a named domain.
\end{abstract}

\section*{Scope}

Formal statements, proof sketches, application records and evidence-status conventions specify the typed ledger. It covers accounting, positivity, barrier, depletion and aggregation results; the interface between a closed ledger and an institutional delay model; the groundwater, phosphate and fisheries records; and operational weak- and strong-sustainability regimes. The certification procedures and reproduction protocol are reported by Abaee (2026a), the accounting-standards analysis by Abaee (2026b), and the institutional delay-dynamics analysis by Abaee (2026c).

The application quantities are bounded descriptive records. The groundwater index is record-relative; the phosphate quotient is an economic reserve-class ratio; and the fisheries quantity is a removals-only pressure time. None is presented here as a calibrated forecast.

\section*{S0. Notation and the declaration discipline}

A typed ledger has state $x(t)\in\Rpos^n$, non-negative primitive fluxes $v(t)\in\Rpos^J$, a declared boundary schedule $b(t)$, a type structure $\mathcal T$, and an incidence operator $S_{\mathcal T}$. Its balance is
\begin{equation}
\dot x=S_{\mathcal T}v+b.\label{eq:s0-balance}
\end{equation}
A composition matrix $C$ maps compartments to declared conserved substance classes, $S=Cx$. A service is a readout $O(x,u,\theta)$, not an additional conserved mass compartment. A margin or barrier is represented by an affine map $m(x)=Gx+a\geq0$ when that form is appropriate.

A primitive column is a \emph{transfer} when its non-zero entries are signed unit entries within one declared type and unit class. It is a \emph{conversion} when the source and product entries carry a declared stoichiometric coefficient and a routed destination. A cross-type sum is not admitted without that conversion declaration. A donor-limited primitive outflow satisfies
\begin{equation}
v_j(t)=0\quad\text{whenever its donor compartment is empty}.\label{eq:s0-donor}
\end{equation}
Every yield below one must route its residual to a represented compartment or an explicitly declared boundary term. These are declaration rules; they are not inferred from a diagnostic label such as ``reserve'', ``waste'' or ``unsustainable portion''.

The status vocabulary is four-valued. \est{} means that the declared obligation has a witness; \nesta{} means that the obligation is stated but not discharged by the declared record; \na{} means that the object has no live obligation of that kind; and \texttt{refuted} requires a witness of the negation. A missing trajectory must not be reported as a refutation.

\section*{S1. Extended theorem inventory}

Table~\ref{tab:inventory} gives the extended inventory. The result labels identify the statements below. Each statement is conditional on the typed ledger, constitutive laws and boundary assumptions stated below; proof sketches identify the assumptions used.

The certification stack is deliberately layered:
\begin{equation}
\begin{gathered}
\texttt{typed}\;\longrightarrow\;\texttt{balanced}\;\longrightarrow\;\texttt{conserved}\;\longrightarrow\;\texttt{admissible}\;\longrightarrow\;\texttt{safe},\\[-2pt]
\texttt{safe}\;\longrightarrow\;\texttt{closure-certified}.
\end{gathered}
\label{eq:certification-stack}
\end{equation}
The arrows name proof obligations, not automatic logical implications. ``Typed'' checks units and routes; ``balanced'' checks the displayed equations; ``conserved'' checks the declared moiety law; ``admissible'' checks non-negativity and donor limits; ``safe'' checks componentwise barriers on a horizon; and ``closure-certified'' additionally requires an admissible return/substitution witness for the declared service. A positive aggregate cannot skip from a lower layer to a higher one.

\begin{longtable}{@{}P{0.19\linewidth}P{0.47\linewidth}P{0.27\linewidth}@{}}
\caption{Extended theorem inventory and scope.}\label{tab:inventory}\\
\toprule
Result & Statement & Boundary of the result \\
\midrule
\endfirsthead
\toprule
Result & Statement & Boundary of the result \\
\midrule
\endhead
\bottomrule
\endfoot
S1.1 Flux reconstruction & Integrating \eqref{eq:s0-balance} reconstructs $x(t)$ from the initial state, typed fluxes and declared boundary terms. & An identity for the declared ledger; it does not identify unmeasured fluxes without additional data or a programme. \\
S1.2 Conservation reduction & If $\ell^{\top}CS_{\mathcal T}=0$, then $\ell^{\top}Cx(T)=\ell^{\top}Cx(0)+\int_0^T\ell^{\top}Cb(t)\dd t$. & Boundary and residual terms remain visible; an undeclared residual is not silently zero. \\
S1.2a Double-count guard & An aggregate row and the disjoint rows that sum to it are mutually exclusive inputs to a total; a derived readout is not added again as a primitive flow. & A data-selection audit is required; overlapping rows can change totals and ratios. \\
S1.3 Flux envelope & Componentwise flux boxes produce lower and upper moiety envelopes by positive/negative matrix parts. & A box is conservative and may contain jointly unrealizable flux choices; it is not a forecast. \\
S1.4 Barrier corollary & If the envelope lies between declared barriers for the full horizon, every trajectory in the box is barrier-safe. & Safety is conditional on the box, the horizon and the declared barriers. \\
S1.5 Uniform-drift bound & If $\dot S\leq-\varepsilon<0$ above a barrier $B$, then $\tau_B\leq(S(0)-B)/\varepsilon$. & The uniform margin is the substantive assumption; positive extraction alone is insufficient. \\
S1.6 Critical-margin budget & A non-negative multiplier satisfying $\lambda^{\top}GS_{\mathcal T}+c^{\top}\leq0$ bounds cumulative service by initial margin plus declared boundary support. & $V>0$ is not a componentwise certificate, and multiplier infeasibility is not proof of safety. \\
S1.7 Closed finite-donor mass & For the closed natural block, $M=N+A_{\mathrm{act}}+A_{\mathrm{geo}}+U$ obeys $\dot M=-qEN-C^{A,\rm lim}$. & The identity is for the declared routing; internal transfers do not exhaust a conserved moiety. \\
S1.8 Orthant invariance & Donor limitation gives forward invariance of $\Rpos^n$ by a face-by-face tangent-cone argument. & It proves non-negativity, not barrier safety, closure, or sustainability. \\
S1.9 No interior rest & With positive extraction effort and the stated constitutive laws, the closed finite-donor block has no interior positive-flux rest point. & The claim is for the closed block and its hypotheses, not for an open target-relaxation completion. \\
S1.10 Rest-set and budget results & The zero-extraction rest set includes extinction, carrying-capacity and frozen-support faces; extraction is $L^1$ over a finite donor budget. & These are structural results, not empirical identification of groundwater, phosphate or fisheries. \\
S1.11 No weighted certification & A non-negative weighted sum does not imply componentwise adequacy on a feasible domain containing a positive aggregate with a negative component. & An additional domain theorem could change the implication; absent one, scalar communication is not certification. \\
S1.12 Depletion is compartmental & A finite crossing concerns a declared compartment, barrier or readout, not total conserved mass. & A reserve-class quotient and a record-relative anomaly index are not automatically physical hitting times. \\
S1.13 First-passage surrogates & Brownian and geometric-Brownian first-passage formulae are valid for their declared surrogate processes. & The surrogates do not stochastically complete the typed ledger and do not identify a physical failure threshold. \\
S1.14 Delay-interface non-reduction & The closed ledger and open institutional delay system share a decline identity under a specialization but do not share a dynamic reduction. & Shared diagnostics do not transfer Hopf, equilibrium or periodic-orbit results. \\
\end{longtable}

\subsection*{S1.1 Balance, conservation and flux envelopes}

Integrating \eqref{eq:s0-balance} gives the reconstruction identity
\begin{equation}
x(T)-x(0)=\int_0^T S_{\mathcal T}v(t)\dd t+\int_0^T b(t)\dd t.\label{eq:reconstruction}
\end{equation}
If $\ell^{\top}CS_{\mathcal T}=0$, multiplication by $\ell^{\top}C$ gives
\begin{equation}
\ell^{\top}Cx(T)=\ell^{\top}Cx(0)+\int_0^T\ell^{\top}Cb(t)\dd t.\label{eq:conservation}
\end{equation}
With a declared residual $d_x$, the right-hand side includes $\int_0^T\ell^{\top}Cd_x(t)\dd t$ and its magnitude is bounded only by the declared residual budget $\varepsilon_\ell(T)$. If that budget is missing, the conservation status is \nesta{}, not established at zero.

\textbf{Proposition S1.2a (Double-counting safeguard).} Let $q_0$ be an aggregate record and let $q_1,\ldots,q_k$ be mutually exclusive component records with $q_0=\sum_iq_i$. A total must use either $q_0$ or $\sum_iq_i$, not $q_0+\sum_iq_i$. The same selector rule applies when a derived readout is computed from a primitive flow: the readout is not added as a second flow.

\emph{Proof.} The mixed total is $q_0+\sum_iq_i=2q_0$ whenever the components exhaust the aggregate and $q_0>0$. For a ratio, the numerator and denominator can be inflated by different overlapping scopes, so the quotient need not be invariant. Choosing one mutually exclusive aggregation level and recording the row selector prevents the duplication. The safeguard is bookkeeping, not a substantive conservation theorem.

Suppose $v(t)\in[\underline v(t),\overline v(t)]$ and $b(t)\in[\underline b(t),\overline b(t)]$.  Write $A^+=\max(A,0)$ and $A^-\!=\max(-A,0)$ entrywise. For moiety row $m$ define
\begin{align}
\varphi_m(t)&=(CS_{\mathcal T})_m^+\underline v(t)-(CS_{\mathcal T})_m^-\overline v(t)+C_m^+\underline b(t)-C_m^-\overline b(t),\\
\psi_m(t)&=(CS_{\mathcal T})_m^+\overline v(t)-(CS_{\mathcal T})_m^-\underline v(t)+C_m^+\overline b(t)-C_m^-\underline b(t),\\
\underline S_m(t)&=S_m(0)+\int_0^t\varphi_m(\tau)\dd\tau,\qquad
\overline S_m(t)=S_m(0)+\int_0^t\psi_m(\tau)\dd\tau.
\end{align}
Then $\underline S_m(t)\leq S_m(t)\leq\overline S_m(t)$. The proof is pointwise: the positive and negative coefficient parts attain the extrema of the linear form on the declared box, and integration preserves the inequality. If $\underline S_m(t)\geq\underline B_m(t)$ and $\overline S_m(t)\leq\overline B_m(t)$ for every $m,t$, the box is a barrier certificate. The certificate is conservative whenever the box includes flux combinations that the coupled dynamics cannot jointly realize.

\subsection*{S1.2 The closed finite-donor block}

Let
\begin{equation}
s=\frac{A_{\mathrm{act}}}{A_{\mathrm{act}}+A_0},\qquad
\sigma=\frac{A_{\mathrm{geo}}}{A_{\mathrm{geo}}+A_{g0}},\qquad
R=rN\left(1-\frac NK\right)s,\qquad T=\kappa_A Ns,
\end{equation}
and $B=R+T$. The donor primitives are
\begin{equation}
e_{GA}=\omega_A A^{\rm eq,intrinsic}\sigma,
\quad e_{AG}=\omega_A A_{\mathrm{act}},
\quad C^{A,\rm lim}=C^A\sigma,
\quad \text{detritus return}=\gamma_UU.
\end{equation}
Under the institutional-failure specialization, the closed block is
\begin{align}
\dot N&=R-qEN,\\
\dot A_{\mathrm{act}}&=-B+e_{GA}-e_{AG}+\gamma_UU,\\
\dot A_{\mathrm{geo}}&=-e_{GA}+e_{AG},\\
\dot U&=T-\gamma_UU.
\end{align}
Summing gives
\begin{equation}
\dot M=-qEN-C^{A,\rm lim},\qquad M=N+A_{\mathrm{act}}+A_{\mathrm{geo}}+U.\label{eq:massblock}
\end{equation}
The identity separates internal routing from block-boundary export. If $C^A=0$, the only block loss is harvest; if harvest is routed partly to a product or detritus compartment, the declared block boundary changes accordingly.

At a coordinate face, every negative primitive outflow has the corresponding donor as a factor or is explicitly donor-limited. Thus the vector field lies in the tangent cone of $\Rpos^4$ and the closed orthant is forward invariant. At a positive-effort interior rest, $\dot U=0$ and the sum of the active and geological equations forces the regeneration/extraction balance; the living-stock equation requires the opposite sign under positive extraction. The two conditions cannot hold at an interior positive-flux rest under the stated constitutive laws. With extraction integrated over $[0,T]$, \eqref{eq:massblock} yields
\begin{equation}
\int_0^T\bigl(qE(t)N(t)+C^{A,\rm lim}(t)\bigr)\dd t\leq M(0).
\end{equation}
This is a finite-budget statement, not a forecast of when a named public indicator crosses a threshold.

\subsection*{S1.3 Critical margins and the aggregation obstruction}

Let $m(x)=Gx+a\geq0$, delivered service $y=c^{\top}v$ with $c\geq0$, and suppose $\lambda\geq0$ satisfies
\begin{equation}
\lambda^{\top}GS_{\mathcal T}+c^{\top}\leq0
\end{equation}
componentwise. Then $V=\lambda^{\top}m(x)$ satisfies
\begin{equation}
\dot V\leq-y+\lambda^{\top}Gb,
\qquad
\int_0^T y(t)\dd t\leq V(x(0))+\int_0^T\lambda^{\top}Gb(t)\dd t.
\end{equation}
The proof is multiplication of the balance by $\lambda^{\top}G$, followed by $v\geq0$ and $V(T)\geq0$. The bound limits the service that can coexist with the declared componentwise margins; it does not make $V$ a replacement for those margins.

For the aggregation obstruction, let $\mathcal B$ be a feasible balance domain. If some $b\in\mathcal B$ satisfies $b_i<0$ and $w^{\top}b>0$ for $w\geq0$, then $w^{\top}b\geq0$ does not imply $b\geq0$. The proof is the exhibited counterexample. For $w=(1/2,1/2)$, $b=(-2,3)$ gives $w^{\top}b=1/2$ while $b_1<0$. An aggregate can communicate or rank, but it cannot certify the conjunction of componentwise barriers without an additional implication proved from the physical domain.

\subsection*{S1.4 Depletion is compartmental and first passage is conditional}

For a stock or readout $S$ and barrier $B$, if $\dot S\leq-\varepsilon<0$ above $B$, then $\tau_B\leq(S(0)-B)/\varepsilon$. The assumption is essential: $\dot S=-kS$ has $S(t)>0$ for every finite $t$, so the zero barrier is never hit in finite time although a positive barrier is reached at $k^{-1}\log(S(0)/B)$. Internal transfers do not exhaust a conserved moiety. A finite zero or threshold crossing therefore refers to a compartment, a functional readout, an accessibility barrier or an economic classification, not automatically to total mass.

For the Brownian surrogate $A(t)=A_0+\mu t+\varsigma W_t$ with $\mu<0$ and record-relative lower barrier $A_{\min}$, the first passage is inverse Gaussian with
\begin{equation}
\nu=\frac{A_0-A_{\min}}{|\mu|},\qquad
\lambda=\frac{(A_0-A_{\min})^2}{\varsigma^2},
\qquad
\mathbb E[T]=\nu,
\qquad
\operatorname{Var}(T)=\frac{\nu^3}{\lambda}.
\end{equation}
For the geometric-Brownian comparison $\dd B_t=-hB_t\dd t+\varsigma B_t\dd W_t$, It\^o's formula gives a logarithmic drift $h+\varsigma^2/2$ and therefore a corresponding inverse-Gaussian mean. These are conditional statements about declared surrogates. They do not complete the mass ledger stochastically.

\section*{S2. Interface with institutional delay dynamics}

\subsection*{S2.1 Exact shared object}

The closed material ledger owns the primitive stock--flow equations, typed routing, conservation, positivity, componentwise margins and finite-donor budget. The institutional delay-dynamics study (Abaee, 2026c) analyses an open institutional completion with memory and review delay. The exact object shared under the single-resource specialization is the decline-pressure identity
\begin{equation}
D(t)=qE(t)N(t)-R\bigl(N(t),A(t)\bigr)=-\dot N(t),
\qquad
\Lambda(t)=[D(t)]_+=[-\dot N(t)]_+.
\label{eq:delay-shared}
\end{equation}
The specialization sets the product, waste and price-response channels and mining intensity to zero, and uses the local stock equation $\dot N=R-qEN$. Equation~\eqref{eq:delay-shared} is an interface identity, not a reduction theorem.

A generic retarded institutional completion can be written abstractly as
\begin{align}
\dot z(t)&=F\bigl(z(t),z_t;\vartheta\bigr),
\qquad z=(N,A_{\mathrm{act}},A_{\mathrm{geo}},U,Z,E),\\
\dot Z(t)&=\mathcal F\bigl(Z_t,\Lambda(t);\vartheta\bigr),
\qquad E(t)=\mathcal E\bigl(Z_t;\vartheta\bigr),
\end{align}
where $z_t(\xi)=z(t+\xi)$ for $\xi\in[-\tau,0]$. This notation identifies the memory input, review delay and effort readout without silently replacing the calibrated equations of Abaee (2026c) by a new model. The delay equations belong to the institutional model; the ledger contributes the typed decline input and its boundary conditions.

\subsection*{S2.2 The two completions}

The closed block uses the intrinsic donor-limited target $A^{\rm eq,intrinsic}$ and
\begin{equation}
e_{GA}=\omega_A A^{\rm eq,intrinsic}\sigma.
\end{equation}
The open institutional completion uses the derived target
\begin{equation}
A^{\rm eq,W}=A^{\rm eq,intrinsic}+\frac{\kappa_AK}{\omega_A},
\end{equation}
which treats continuing donor support as an imposed completion. The interface calculation uses $A^{\rm eq,intrinsic}=50$, $A_{\mathrm{act},*}=397.87$ and $A^{\rm eq,W}=5050$.

At the same state, the active-pool vector fields differ by
\begin{equation}
\omega_A\bigl(A^{\rm eq,W}-A^{\rm eq,intrinsic}\sigma\bigr)-\gamma_UU
=\kappa_AK-\gamma_UU
\end{equation}
under the declared scale separation $\sigma\approx1$. At the working point, the difference is approximately $0.535$ stock units per year; at $U=0$ the scale is $\kappa_AK=5.000$. The difference is structural: the institutional completion omits the detritus return carried by the closed field and uses a derived target.

\subsection*{S2.3 Non-reduction result}

\textbf{Proposition S2.1 (One-way projection).} Under the institutional-failure specialization, if the right-hand sides for $(N,A_{\mathrm{act}},A_{\mathrm{geo}},U,Z,E)$ contain no excluded macroeconomic state, that block is an exact projection for every parameter value. The converse projection is not available when excluded variables are driven by the block.

\emph{Proof sketch.} A vector field involving only the displayed block defines an invariant projection; no singular limit or timescale separation is needed. The statement concerns dependency structure only and does not imply that the projected block is the closed physical ledger.

\textbf{Theorem S2.2 (Non-reduction of the open institutional completion).} There is no exact dynamic reduction, regular perturbation or finite-time tracking correspondence from the closed primitive finite-donor ledger to the open institutional completion under the declared definitions.

\emph{Reason.} Five independent obstructions are recorded:
\begin{enumerate}
\item the intrinsic and derived targets differ by $\kappa_AK/\omega_A=5000$ stock units;
\item the two active-pool fields differ at the same state by $\kappa_AK-\gamma_UU$, an $O(1)$ to $O(\kappa_AK)$ term;
\item the working point requires continuing geological support and is not a rest point of the closed finite-donor system;
\item $\varepsilon_G(T)=G_0^{-1}\int_0^T|e_{GA}-e_{AG}|\dd t$ is a draw diagnostic, not a trajectory-tracking error; and
\item the closed ledger makes sustained extraction $L^1$ in time, preventing indefinite positive-flux persistence on a finite donor budget.
\end{enumerate}

The permitted relation is analogy for shared mechanism language plus reconstruction of omitted flows. Hopf crossings, periodic orbits and stability thresholds of the frozen-donor institutional system do not transfer to the closed ledger. Conversely, the open institutional system can be used as a diagnostic completion whose omitted turnover and imposed recharge are named, not as a theorem-preserving reduction.

\subsection*{S2.4 Frozen-donor limit and finite-budget reading}

With $G=G_0g$ and $g(0)=1$,
\begin{equation}
\dot g=-G_0^{-1}(e_{GA}-e_{AG}).
\end{equation}
The limit $G_0\to\infty$ freezes the donor coordinate but does not restore the institutional completion's derived target; it is therefore not a regular perturbation of the working vector field. If a sustained extraction flux satisfies $c(t)\geq c_0>0$ over an interval, the closed donor budget gives the order bound $T\leq G_0/c_0$. This is a budget bound, not an asymptotic estimate of delay-induced cycles. Whether a frozen-donor Hopf structure persists as a slowly drifting transient is an open slow-passage question.

\section*{S3. Full application tables and status records}

\subsection*{S3.1 Classification matrix}

\small
\begin{longtable}{@{}P{0.18\linewidth}P{0.16\linewidth}P{0.16\linewidth}P{0.16\linewidth}P{0.16\linewidth}@{}}
\caption{Three records and the time-like quantities they do or do not instantiate.}\label{tab:classification}\\
\toprule
Quantity & Question & G3P & Phosphate & Fisheries \\
\midrule
\endfirsthead
\toprule
Quantity & Question & G3P & Phosphate & Fisheries \\
\midrule
\endhead
\bottomrule
\endfoot
$J_A^{\rm gross},H_A^{\rm gross}$ & turnover/\allowbreak{}dependency & no & no & gross-loss analogue only \\
$H_A^{\rm loc}$ & frozen net-rate ratio & no: anomaly, not stock & no: classification, not stock & no: no net $\dot B$ \\
$T_A$ & scenario hitting time & no & no & no \\
Object status & strongest permissible reading & record-relative statistical index & economic reserve-class quotient & removals-only pressure scale \\
\end{longtable}
\normalsize

\subsection*{S3.2 Application contract: observable, denominator and non-claim}

The numeric tables are most safely read with the following contract. ``Barrier'' means the declared reference used by the arithmetic; it does not mean that the source has identified a physical failure floor.

\scriptsize
\begin{longtable}{@{}P{0.10\linewidth}P{0.11\linewidth}P{0.085\linewidth}P{0.10\linewidth}P{0.11\linewidth}P{0.13\linewidth}P{0.10\linewidth}P{0.13\linewidth}@{}}
\caption{Application contract for the three public-data records.}\label{tab:application-contract}\\
\toprule
Record & Observable & Unit & Denom. & Barrier/\allowbreak{}reference & Interpretation & Source/\allowbreak{}status & Explicit\allowbreak{} non-claim \\
\midrule
\endfirsthead
\toprule
Record & Observable & Unit & Denom. & Barrier/\allowbreak{}reference & Interpretation & Source/\allowbreak{}status & Explicit\allowbreak{} non-claim \\
\midrule
\endhead
\bottomrule
\endfoot
G3P/\allowbreak{}GRACE groundwater & basin-mean anomaly $a$ and fitted trend $\widehat{\dot a}$ & cm; cm/yr & $[-\widehat{\dot a}]_+$ & record-window minimum & persistence of an observed anomaly relative to its own window & G3P v1.12; GRACE-\allowbreak{}derived anomaly; basin extraction caveat & not absolute aquifer volume, recharge rate, physical depletion or forecast \\
USGS phosphate & reserve mass $G_{\rm reserve}$ and production $C_G$ & kt; kt/yr & $C_G$ & reserve classification & arithmetic reserve-life ratio & USGS MCS 2026 world pin; country rows retain vintage notes & not geological exhaustion; reserves are not resources and classification can change \\
RAM Legacy fisheries & $\mathrm{SSB}_{\rm now}$, $B_{\rm lim}$ and $F_{\rm now}$ & biomass; yr$^{-1}$ & $F_{\rm now}$ & $B_{\rm lim}=0.2\max\mathrm{SSB}$ & removals-only pressure time & archived 43-stock pull plus v4.44/v4.66 public comparisons & not net biomass decline, demographic hitting time, policy forecast or stock-recruit model \\
\end{longtable}
\normalsize

\subsection*{S3.3 G3P groundwater anomaly table}

For a basin anomaly $a(t)$, the reported index is
\begin{equation}
L_{\rm hist}^{\rm anom}=
\frac{a_{\rm latest}-a_{\rm hist,min}}{[-\widehat{\dot a}]_+}.
\end{equation}
The product supplies anomalies relative to a reference period, not absolute aquifer volumes. The window, basin mask, anomaly reference and fitted trend are part of the record status.

\begin{longtable}{@{}P{0.27\linewidth}P{0.13\linewidth}P{0.15\linewidth}P{0.17\linewidth}P{0.18\linewidth}@{}}
\caption{G3P basin-row extraction used to exhibit the record-relative index.}\label{tab:g3p}\\
\toprule
Basin & Trend (cm/yr) & 2023 anomaly (cm) & Implied window minimum (cm) & Index (yr) \\
\midrule
\endfirsthead
\toprule
Basin & Trend (cm/yr) & 2023 anomaly (cm) & Implied window minimum (cm) & Index (yr) \\
\midrule
\endhead
\bottomrule
\endfoot
Indo-Gangetic (N. India)\textsuperscript{a} & $-49.7$ & $-414$ & $-548$ & $\approx2.7$ \\
North China Plain & $-18.6$ & $-145$ & $-292$ & $\approx7.9$ \\
Central Valley (US) & $-16.1$ & $-84$ & $-237$ & $\approx9.5$ \\
La Mancha (Spain) & $-3.2$ & $-20$ & $-88.5$ & $\approx21.4$ \\
High Plains (US) & $-7.9$ & $-160$ & $-160$ & already at minimum \\
global mean & $-0.4$ & $-14$ & $-33.0$ & $\approx47.5$ \\
\end{longtable}

\textsuperscript{a} The Indo-Gangetic magnitude is quarantined. It sits well beyond published basin-mean trends and must be re-derived from the product's basin masks before numerical reuse. The classification status does not depend on the magnitude.

\subsection*{S3.4 Phosphate reserve table}

The arithmetic reserve-life quotient is $T_{\rm reserve}=G_{\rm reserve}/C_G$. On the pinned record it is approximately $309$ years for world reserves. It is an economic classification quotient, not a geological exhaustion forecast.

\begin{longtable}{@{}P{0.25\linewidth}P{0.22\linewidth}P{0.24\linewidth}P{0.22\linewidth}@{}}
\caption{Phosphate reserve-class table and arithmetic quotient.}\label{tab:phosphate}\\
\toprule
Country/class & Reserves (kt) & Reserve-life horizon (yr) & Implied production (kt/yr) \\
\midrule
\endfirsthead
\toprule
Country/class & Reserves (kt) & Reserve-life horizon (yr) & Implied production (kt/yr) \\
\midrule
\endhead
\bottomrule
\endfoot
China & $3{,}400{,}000$ & $\approx28$ & $121{,}429$ \\
United States & $1{,}000{,}000$ & $\approx45$ & $22{,}222$ \\
Jordan & $820{,}000$ & $\approx62$ & $13{,}226$ \\
Morocco & $50{,}000{,}000$ & $\approx1{,}250$ & $40{,}000$ \\
Australia\textsuperscript{a} & $5{,}800{,}000$ & $\approx2{,}088$ & $2{,}778$ \\
World reserves & $74{,}000{,}000$ & $\approx309$ & $239{,}482$ \\
World resources, $\varepsilon=0.10$ & $>300{,}000{,}000$ & $>1{,}125$ & $240{,}000$ \\
\end{longtable}

\textsuperscript{a} The Australia row carries a pre-2026 reserve vintage and is not used to establish the world classification. Reserve membership changes with prices, technology, exploration and regulation; a resources threshold is a different declared object.

\subsection*{S3.5 Fisheries removals-only table}

For the declared comparison process $\dot B=-F_{\rm now}B$, with $B_{\rm lim}=0.2\max\mathrm{SSB}$,
\begin{equation}
\Theta_F=\frac{\log(\mathrm{SSB}_{\rm now}/B_{\rm lim})}{F_{\rm now}}.
\end{equation}
It omits recruitment, growth, maturation, natural mortality, density dependence, environmental forcing and future policy.

\begin{longtable}{@{}P{0.25\linewidth}P{0.16\linewidth}P{0.18\linewidth}P{0.24\linewidth}@{}}
\caption{Fisheries cohort and release sensitivity record.}\label{tab:fisheries}\\
\toprule
Record & Stocks & Median $\Theta_F$ (yr) & Status and caveat \\
\midrule
\endfirsthead
\toprule
Record & Stocks & Median $\Theta_F$ (yr) & Status and caveat \\
\midrule
\endhead
\bottomrule
\endfoot
Archived selected cohort & 43 & $\approx1.8$ & Selected cohort; eight zero entries at or below the reference; distinct from both public releases \\
Archived positive sub-cohort & 35 & $\approx2.9$ & $F>0$ and $\mathrm{SSB}_{\rm now}>B_{\rm lim}$; still removals-only \\
RAM Legacy v4.44 public release & 415 & $\approx2.57$ & Same protocol; public-release comparison \\
RAM Legacy v4.66 public release & 454 & $\approx3.39$ & Same protocol; broad-cohort result \\
\end{longtable}

The selected 43-stock record is a class-specific diagnostic rather than a statistic of all assessed fisheries. The broad-release result characterizes database-cohort dependence and is not an uncertainty interval around a demographic forecast.

\section*{S4. Operational weak and strong sustainability regimes}

\subsection*{S4.1 Declared closure objects}

Let $\mathcal V(x,t)$ be the set of non-negative primitive flux vectors satisfying the typed incidence rules, donor limits, conversion declarations and boundary constraints at $(x,t)$. Let $D(t)$ be a declared service-demand vector and $P$ map return fluxes to the uses they cover. Let $m(x)=Gx+a$ be the declared componentwise margin. Define the closure capacity
\begin{equation}
\Lambda^*(x,t)=\sup\left\{\lambda\geq0:\ \exists v\in\mathcal V(x,t),\quad
\lambda D(t)\leq Pv,\quad m(x)\geq0\right\}.
\label{eq:closure-capacity}
\end{equation}
The inequality in \eqref{eq:closure-capacity} is a declared coverage relation; it does not equate a service unit with a material unit. Cross-type substitution enters only through a declared conversion column and its coefficient. An external input is a boundary input, not internal regeneration.

For a horizon $T$, let $\mathcal A_T$ be the set of admissible trajectories satisfying the typed balance, donor and conversion rules, declared boundary schedule and margin constraints on $[0,T]$. The closure object is therefore a feasible-set statement, not a scalar sustainability score.

\subsection*{S4.2 Definitions}

\textbf{Definition S4.1 (Operational weak-closure regime).} The declared ledger is in the weak-closure regime on $[0,T]$ if there exists a trajectory in $\mathcal A_T$ that meets the declared demand and whose return network has $\Lambda^*(x(t),t)\geq1$ at the relevant use times. Declared substitutions are admissible only when their type conversion, coefficient, destination and residual route are present.

\textbf{Definition S4.2 (Operational strong-component regime).} The declared ledger is in the strong-component regime on $[0,T]$ when a critical component or barrier is a live certification obligation and either (i) no admissible conversion/return route covers its demand at the declared horizon, $\Lambda^*<1$, or (ii) a componentwise margin is binding or violated. A surplus in another component does not clear that obligation.

These definitions are operational readings of a typed ledger. They do not assert that all ethical, distributive or intergenerational meanings of weak and strong sustainability reduce to $\Lambda^*$ and $m$.

\subsection*{S4.3 Conditional regime proposition}

\textbf{Proposition S4.3 (What the regime labels certify).} If the weak-closure conditions hold, the declared return network can cover the declared demand on the stated horizon while satisfying the declared componentwise margins. If the strong-component condition holds, the record identifies a binding component, a deficit in declared return capacity, or an unestablished route; no non-negative scalar aggregate certifies the missing component without an additional domain implication.

\emph{Proof.} The first statement is the feasibility witness in \eqref{eq:closure-capacity} together with membership in $\mathcal A_T$. For the second, either $\Lambda^*<1$ by definition or a live componentwise margin fails. A scalar aggregate that remains positive while a component is negative is not equivalent to the conjunction of component inequalities; the counterexample $b=(-2,3)$ with $w=(1/2,1/2)$ supplies the witness. \(\square\)

\subsection*{S4.4 Formalization boundary}

A universal weak/strong-sustainability theorem would require a complete class of closure maps, admissible substitutions, boundary flows, time horizons, critical-stock definitions and institutional valuation rules. These definitions therefore apply only under the stated physical substitution and closure assumptions; cross-type compensation is not a componentwise certificate.

\begin{longtable}{@{}P{0.24\linewidth}P{0.30\linewidth}P{0.34\linewidth}@{}}
\caption{Reporting template for the two operational regimes.}\label{tab:regimes}\\
\toprule
Field & Weak-closure reading & Strong-component reading \\
\midrule
\endfirsthead
\toprule
Field & Weak-closure reading & Strong-component reading \\
\midrule
\endhead
\bottomrule
\endfoot
Physical route & Declared conversions and return routes cover demand & Route absent, insufficient or not identified \\
Closure capacity & $\Lambda^*\geq1$ at the declared use times & $\Lambda^*<1$ or not established \\
Component margin & $m_j(t)\geq0$ on the declared horizon & A critical $m_j$ binds, fails or is not established \\
Scalar aggregate & Communication or ranking only & Cannot clear the component obligation \\
Status wording & State the witness and horizon & Name the deficit, missing route or binding barrier \\
\end{longtable}

\section*{S5. Evidence status and reproducibility}

The application tables report source-vintage arithmetic and descriptive diagnostics. The methods and certification procedures are specified by Abaee (2026a); the accounting-standards analysis is given by Abaee (2026b); and the delay-dynamics model is given by Abaee (2026c). The institutional interface stated here does not transfer the delay study's Hopf or periodic-orbit results to the closed ledger. The stochastic first-passage processes are declared surrogates and do not constitute stochastic completions of the mass ledger.


\section*{References}

Ahuja, R.K., Magnanti, T.L., Orlin, J.B., 1993. \emph{Network Flows: Theory, Algorithms, and Applications}. Prentice-Hall.

Aubin, J.-P., 1991. \emph{Viability Theory}. Birkh\"auser.

Brunner, P.H., Rechberger, H., 2004. \emph{Practical Handbook of Material Flow Analysis}. Lewis Publishers.

Chhikara, R.S., Folks, J.L., 1989. \emph{The Inverse Gaussian Distribution}. Marcel Dekker.

Daly, H.E., 1990. Toward some operational principles of sustainable development. \emph{Ecological Economics} 2, 1--6.

Ekins, P., Simon, S., Deutsch, L., Folke, C., De Groot, R., 2003. A framework for the practical application of the concepts of critical natural capital and strong sustainability. \emph{Ecological Economics} 44, 165--185.

Feinberg, M., 2019. \emph{Foundations of Chemical Reaction Network Theory}. Springer.

G\"untner, A., Sharifi, E., Haas, J., et al., 2024. Global Gravity-based Groundwater Product (G3P), v1.12. GFZ Data Services. \url{https://doi.org/10.5880/G3P.2024.001}.

Jacquez, J.A., Simon, C.P., 1993. Qualitative theory of compartmental systems. \emph{SIAM Review} 35, 43--79.

Munda, G., Nardo, M., 2009. Noncompensatory/nonlinear composite indicators for ranking countries: a defensible setting. \emph{Applied Economics} 41, 1513--1523.

Neumayer, E., 2013. \emph{Weak versus Strong Sustainability: Exploring the Limits of Two Opposing Paradigms}, 4th ed. Edward Elgar.

Oksendal, B., 2003. \emph{Stochastic Differential Equations}, 6th ed. Springer.

Redner, S., 2001. \emph{A Guide to First-Passage Processes}. Cambridge University Press.

Ricard, D., Minto, C., Jensen, O.P., Baum, J.K., 2012. Examining the knowledge base and status of commercially exploited marine species with the RAM Legacy Stock Assessment Database. \emph{Fish and Fisheries} 13, 380--398.

Tilton, J.E., 2003. \emph{On Borrowed Time? Assessing the Threat of Mineral Depletion}. Resources for the Future.

U.S. Geological Survey, 2026. \emph{Mineral Commodity Summaries 2026: Phosphate Rock}. USGS.

Abaee, A., 2026a. \emph{Certifying a Typed Ledger: The Predicates, the Programmes, the Vintages, and the Reproduction Bundle}. Figshare preprint. \url{https://doi.org/10.6084/m9.figshare.33942469}

Abaee, A., 2026b. \emph{What the Accounts Settle, and What They Leave Open: Depletion as a Cost of Production, and the Missing Step from a Rate to a Horizon}. Figshare preprint. \url{https://doi.org/10.6084/m9.figshare.33942487}

Abaee, A., 2026c. \emph{Delay-Induced Regime Change in Harvested Stocks: The Mobilising and Protective Channels of Institutional Feedback, and the Review Interval as Control}. Zenodo preprint. \url{https://doi.org/10.5281/zenodo.22554217}

\end{document}
